\documentclass[12pt, notitlepage]{article}   	% use "amsart" instead of "article" for AMSLaTeX format
\usepackage{geometry}                		% See geometry.pdf to learn the layout options. There are lots.
\geometry{a4paper}                   		% ... or a4paper or a5paper or ... 
%\geometry{landscape}                		% Activate for rotated page geometry
\usepackage[parfill]{parskip}    		% Activate to begin paragraphs with an empty line rather than an indent
\usepackage{graphicx}				% Use pdf, png, jpg, or eps§ with pdflatex; use eps in DVI mode
								% TeX will automatically convert eps --> pdf in pdflatex

\usepackage{hyperref}
		
%SetFonts

\usepackage[T1]{fontenc}
\usepackage[utf8]{inputenc}

\usepackage{tgbonum}

%SetFonts

\title{
	\textbf{
		BIOS 998-002
	} \\
	\large Ecosystem Modeling \\
	\large Spring 2027
}

\date{\vspace{-5ex}}

\begin{document}

{\fontfamily{phv}\selectfont %select helvetica (code = phv)

\maketitle

\section{Course Description}
Students in this course will learn the core concepts behind systems thinking and model building.
Through course-based model development,
students will learn the skills necessary to develop their own model based on their interests
and to understand
the workings of models developed by others. Models will be developed independently and 
connected collaboratively into a class-scale ecosystem model.
The course is intended for graduate students interested in developing skills in
quantitative theory development, systems thinking, and model building. Students will also
learn best practices in R coding, version control, and large project collaboration.

\subsection{Class Time and Location}
Tuesdays and Thursdays 12:30-13:45

Beadle N172

\subsection{Instructor}
Dr. Nick Smith \par
Manter 402C \par
nsmith87@unl.edu \par

\subsection{Office Hours}
Thursdays 14:00-14:45 \par
Manter 402C \par

\subsection{Recommended Texts}
Climate Change and Terrestrial Ecosystem Modeling
by Bonan
\url{https://doi.org/10.1017/9781107339217}\par

Principles of Terrestrial Ecosystem Ecology (2nd Edition; 2011) 
by Chapin, Matson, and Vitousek \par
The book can be accessed from Springer here: 
\url{https://link.springer.com/book/10.1007/978-1-4419-9504-9}. Click on "Access this title on 
SpringerLink." It can also be accessed through the TTU library.

\section{Course Materials}
All course materials, including lecture slides, readings, activities, and code 
will be posted to a GitHub repository for the course.
The primary repository address is
\url{https://github.com/SmithEcophysLab/ecosystem_modeling_sp2027}.
The repository will include the syllabus and all other miscellaneous class materials as the semester
progresses. A README file will contain information on the repository, including
links to different sections at 
\url{https://github.com/SmithEcophysLab/ecosystem_modeling_sp2027/README.md}.

\section{Learning Objective}
This course will broadly focus on understanding the principles underlying the development of
process-based models of ecosystems. An emphasis will be placed on how to build these
models to understand the interactions between ecosystems and drivers of global
change, including climate change, changes in atmospheric gas concentration, and eutrophication.
Course topics will be taught through the lens of classical and more contemporary 
primary and secondary literature. There will be a strong emphasis on hands-on modeling skill
development. Topics will be flexible and modified to match student interests where possible.

\section{Attendance Policy}
Attendance will not be taken, but is strongly recommended. 
In class activity points will only be granted if students are in class.
Makeups will not be granted.

\section{Course Assessment}
\subsection{\textit{Participation and Engagement}}
Being an active and engaged participant in the class will benefit your understanding
of material as well as your peers'. Examples include asking questions, providing feedback,
and facilitating discussion.

\subsection{\textit{Quizzes}}
Short quizzes will be given periodically to test student knowledge of core concepts and to
stimulate discussion. In some cases, quizzes will be developed and administered
by class weekly co-leads.

\subsection{\textit{Homework Assignments}}
Short homework assignments will be given each week based on the Tuesday lectures and
Thursday model development discussions. These will typically be assigned on Thursday and
due the following Tuesday.

\subsection{\textit{Weekly co-leads}}
Throughout the semester, students will be asked to co-lead on the week's discussion
topic with Dr. Smith. This will consist of leading a jigsaw-style discussion of a literature article on Tuesdays
and developing a quiz for students on Thursdays.
Students will be evaluated on their ability to respond accurately to their peers' questions
as well as their ability to summarize and generate discussion on the week's topic.

\subsection{\textit{Model module development}}
The primary semester project will be to develop a module for the class ecosystem
model based on the student's interest. The module will be written in R and must be able to
run as a stand-alone module as well as in connection with the larger class model.
The student will be required to write a full description of the module to coincide with the
code in both README and manuscript style format.

\section{Grading}
Participation and Engagement: 25\% \par
Quizzes: 5\% \par
Homework Assignments: 20\% \par
Weekly co-lead: 15\% \par
Module proposal: 10\% \par
Module presentation: 5\% \par
Final module: 20\% \par

Grades will be made available on Blackboard. 
All grades posted at the end of the course will be final, 
unless an error has been made in their calculation.
Please contact Dr. Smith if you feel your grade has been calculated incorrectly.

\section{Grading Scale}
A: $\geq$ 90\% \par
B: 80 – 90\% \par
C: 70 – 80\% \par
D: 60 – 70\% \par
F: $\leq$ 59.9\% \par

\section{Missing In-class Activities}
Students will be required to be in class to receive in-class activity points. 
Please contact Dr. Smith if you plan to miss class for a university function 
\textit{prior to class}. If class is missed due to an illness, 
please let Dr. Smith know as soon as possible. Documentation will need to be provided
in order to be able to make up any missed work.

\newpage

\section*{Schedule of Topics by Week}
01/12/2027 – Introductions, semester planning, and goals \par
01/19/2027 – Systems Thinking; building systems diagrams \par
01/26/2027 – Terrestrial Ecosystem cycles; downloading and using R \par
02/02/2027 – Ecosystem Models; GitHub and the class model environment \par
02/09/2027 – 3 R's of Ecosystem Modeling; working collaboratively in GitHub \par
02/16/2027 – Integrating Ecology into Models; writing functions in R \par
02/23/2027 – Eco-evolutionary Optimality Theory; model repository structure \par
03/02/2027 – \textbf{Module proposal presentations} \par
03/09/2027 – \textbf{Module proposal presentations} \par
03/16/2027 – NO CLASS; SPRING BREAK \par
03/23/2027 – Model development: Carbon \par
03/30/2027 – Model Development: Water \par
04/06/2027 - TBD \par
04/13/2027 – Model Development: Nutrients \par
04/20/2027 – \textbf{Module presentations} \par
04/27/2027 – \textbf{Module presentations}; \textbf{Final module due} \par

\section*{General Weekly Schedule}
Generally, each Tuesday will consist of a lecture and jigsaw-style paper discussion. Thursdays
will consist of a quiz, paper discussion, and class model development exercise.

} %end font selection

\end{document} 
